% =============================================================================
% APPENDIX: REPRODUCIBILITY
% =============================================================================
\appendix
\section{Reproducibility Details}
\label{app:repro}
This appendix specifies the exact splits, seeds, model versions, hyperparameters, and code needed to reproduce every reported number, and states the provenance of each value in Table~\ref{tab:comparative}.

\subsection{Datasets and Splits}
\textbf{PlantVillage (controlled, Table~\ref{tab:comparative}).} We use the 11-class novel subset (Apple, Grape, Potato; 9{,}385 images) with frozen embeddings cached once. Reproducing these controlled runs requires only the cached embeddings and the provided script. \emph{11-class regime:} 30 Monte-Carlo episodes; episode $e$ uses \texttt{random.Random(42+17e)}; for each class the image indices are shuffled, the first 5 form the support and all remaining images form the query (9{,}330 queries/episode; support $\cap$ query $=\emptyset$). \emph{5-way $K$-shot regime:} 50 episodes per $K\in\{1,5,10,15\}$; a single \texttt{random.Random(42)} stream samples 5 of the 11 classes per episode (\texttt{rng.sample(range(11),5)}), shuffles each class, takes the first $K$ as support and the rest as query. All methods share identical support/query indices and cached backbone outputs within each episode. Zero-shot BioCLIP and Tip-Adapter use the BioCLIP branches, whereas centroid and Fisher classifiers evaluate joint BioCLIP+DINOv2 representations to isolate the effect of the adaptation rule.

\textbf{PlantDoc (field domain, labeled target-domain support).} The 27 classes shared by both splits are used; the training-only class (Tomato two-spotted spider mites, 2 images) is dropped, leaving 2{,}668 training and 252 test images. 5-way 5-shot: 50 episodes, \texttt{random.Random(42)}, \texttt{rng.sample(range(27),5)}; 5 support leaves/class drawn from the training split, all test images of the sampled classes as queries. 27-way: 135 support leaves (5/class, 8-view augmented to 1{,}080 feature vectors) and 252 field queries. PlantDoc images come from the original benchmark~\cite{singh2020plantdoc}; the main tables use the 2{,}922-image working redistribution described above, and Appendix~\ref{app:plantdoc_official} compares it with the official release.

\subsection{Full Shot Ladder}
\label{app:shot_ladder}
Table~\ref{tab:shot_ladder} reports diagnostic accuracy across 1-shot, 5-shot, 10-shot, and 15-shot regimes (50 episodes/setting; Appendix~\ref{app:repro}). Zero-shot BioCLIP requires no support labels; minor variation reflects query sampling variance. Published baselines provide literature context under their original protocols (e.g., SC-FSL reports 95\% confidence intervals rather than standard deviations, and class pools differ); studies without verifiable 5-shot metrics (e.g., MAFDE-DN4~\cite{zhao2024mafde}) are discussed qualitatively in Section~\ref{sec:related_work}.

\begin{table*}[htbp]
\caption{Full PlantVillage 5-way shot ladder (accuracy, \%, single-view reference features).}
\label{tab:shot_ladder}
\centering
\footnotesize
\setlength{\tabcolsep}{4pt}
\begin{tabularx}{\textwidth}{@{}X c c c c@{}}
\toprule
\textbf{Method} & \textbf{1-shot} & \textbf{5-shot} & \textbf{10-shot} & \textbf{15-shot} \\
\midrule
\multicolumn{5}{@{}l}{\textit{Controlled evaluations, 11-class pool}} \\
Zero-shot BioCLIP & 38.89$\pm$12.78 & 38.86$\pm$12.79 & 38.83$\pm$12.81 & 38.79$\pm$12.83 \\
Tip-Adapter & 45.15$\pm$13.42 & 68.63$\pm$12.80 & 86.01$\pm$7.49 & 93.23$\pm$4.00 \\
Prototypical Networks & 88.44$\pm$9.13 & 97.45$\pm$2.03 & 98.35$\pm$1.45 & 98.57$\pm$1.27 \\
BioDINO-Edge (inductive) & 88.45$\pm$9.12 & 97.96$\pm$1.85 & 99.10$\pm$0.87 & 99.41$\pm$0.60 \\
BioDINO-Edge (transductive) & 93.42$\pm$6.83 & 98.77$\pm$1.68 & 99.41$\pm$0.60 & 99.63$\pm$0.39 \\
\midrule
\multicolumn{5}{@{}l}{\textit{Published results (context only)}} \\
SC-FSL~\cite{mu2024few} & 78.55$\pm$0.81 & 92.90$\pm$0.47 & -- & -- \\
AgriFewNet~\cite{babu2025agrifewnet} & 87.30 & 94.80 & 97.10 & -- \\
VSF-Trans (Swin-T)~\cite{zhou2026vsf} & 79.01 & 90.51 & -- & -- \\
\bottomrule
\end{tabularx}

\vspace{3pt}
\begin{minipage}{\textwidth}
\footnotesize
\textit{Notes.} Controlled results are mean $\pm$ std over 50 episodes per setting. Published baselines are quoted as reported for context under their original protocols; dashes indicate unreported settings.
\end{minipage}
\end{table*}

\subsection{PlantDoc Release Robustness}
\label{app:plantdoc_official}
The main tables evaluate 2{,}668 training and 252 test images from a 2{,}922-image working redistribution of PlantDoc~\cite{singh2020plantdoc} (omitting two images from the unrepresented spider-mite class). To verify partition robustness, we repeated headline evaluations on the official GitHub release (2{,}342 training / 236 test images) under identical seeds, support sizes, 8-view augmentation, and shrinkage parameters. As reported in Table~\ref{tab:plantdoc_official}, performance differences between the working redistribution and official benchmark release are not statistically significant: 27-way transductive accuracy yields 52.38\% (working) vs.\ 56.78\% (official) ($z=0.98, p=0.33$), and 5-way 5-shot episodic accuracy yields 84.94\% (working) vs.\ 82.17\% (official) (Welch $t=1.15, p=0.25$).

An MD5 (Message Digest Algorithm 5) hash audit identified six test images byte-identical to training samples (four with conflicting annotations in the benchmark). Crucially, none belong to the 135 support leaves ($27 \times 5$), ensuring zero influence on prototype or covariance adaptation. Excluding these six test queries bounds the 27-way accuracy shift to $\pm1.24$\,pp. The partition manifest documenting all 2{,}922 records is published in \path{plantdoc_split_manifest.csv}.

\begin{table}[htbp]
\caption{Headline PlantDoc results on the working redistribution and on the official GitHub release (8-view support expansion).}
\label{tab:plantdoc_official}
\centering
\footnotesize
\setlength{\tabcolsep}{4pt}
\begin{tabular*}{\columnwidth}{@{\extracolsep{\fill}}lcc}
\toprule
\textbf{Regime} & \textbf{Working copy} & \textbf{Official release} \\
 & (2{,}668 / 252) & (2{,}342 / 236) \\
\midrule
27-way inductive & 48.81 & 54.66 \\
27-way transductive & 52.38 & 56.78 \\
27-way macro F1 & 52.52 & 57.33 \\
5-way 5-shot inductive & 82.60$\pm$11.49 & 80.33$\pm$13.94 \\
5-way 5-shot transductive & 84.94$\pm$11.15 & 82.17$\pm$12.86 \\
5-way 5-shot macro F1 & 84.55$\pm$11.20 & 81.48$\pm$13.48 \\
\bottomrule
\end{tabular*}

\vspace{3pt}
\begin{minipage}{\columnwidth}
\footnotesize
\textit{Notes.} Both columns evaluate 8-view support expansion under identical random seeds and shrinkage parameters. Differences between distributions are not statistically significant ($p > 0.25$).
\end{minipage}
\end{table}

\subsection{Models and Hyperparameters}
Table~\ref{tab:hyperparameters} summarizes the architecture specifications, hyperparameter constants, and numerical solver configurations across all evaluation regimes. The Python reference implementation uses \texttt{scikit-learn} with \texttt{solver=lsqr}, while the mobile deployment uses dense Cholesky factorization in EJML. Because conjugate gradients in \texttt{lsqr} solve the identical normal equations, both implementations converge within machine precision. Quantization audits (Table~\ref{tab:quantization}) confirm that end-to-end transductive accuracy changes by at most 0.40\,pp on PlantDoc and 0.24\,pp on PlantVillage under INT8 ONNX execution.

\begin{table}[htbp]
\caption{System hyperparameters and implementation specifications.}
\label{tab:hyperparameters}
\centering
\footnotesize
\setlength{\tabcolsep}{4pt}
\begin{tabularx}{\columnwidth}{@{}l X c@{}}
\toprule
\textbf{Component / Parameter} & \textbf{Symbol / Description} & \textbf{Value / Setting} \\
\midrule
DINOv2 backbone & Frozen visual patch encoder ($d_{\text{dino}}$) & ViT-S/14 (384-D) \\
BioCLIP backbone & Frozen botanical vision--language ($d_{\text{bio}}$) & ViT-B/16 (512-D) \\
Joint representation & Concatenated embedding ($\mathbf{z}$) & 896-D \\
Logit scaling & Cosine temperature ($\gamma$) & 20.0 \\
Fisher blend weight & Inductive blend factor ($w_{\text{lda}}$) & 0.75 \\
Uncertainty gate & Shannon-margin threshold ($\tau$) & 0.35 \\
Transductive support weight & Pseudo-label anchor weighting & $3\times$ \\
Anchor buffer capacity & Maximum anchors per class & 20 \\
Anchor confidence gate & Pseudo-label entry threshold & 0.70 \\
Covariance shrinkage & Regularization estimator & Ledoit--Wolf \\
Singular ridge floor ($N=C$) & Degenerate 1-shot limit & Mean feature variance \\
Tip-Adapter baseline & Blending constants ($\alpha, \beta$) & 1.0, 5.5 \\
Mobile / Server solvers & Dense Cholesky (EJML) $\equiv$ Iterative LSQR & Equivalent ($<10^{-6}$) \\
\bottomrule
\end{tabularx}
\end{table}

\subsection{Environment and Code}
Reference environment: Python 3.12.1, torch 2.6.0 (Compute Unified Device Architecture (CUDA) 12.4 build), scikit-learn 1.9.1, numpy 2.5.4, scipy 1.18.1, timm 1.0.30, open\_clip 3.3.0. Feature extraction ran on an NVIDIA RTX 4070 Laptop GPU, while on-device latency, memory, and thermal profiling executed on the ARM64 CPU of the Galaxy S24 FE. The open-source repository (\url{https://github.com/reyhane-code/BioDINO-Edge}) provides all reproduction scripts and models: (i)~core framework and inference in \path{pipeline/biodino_edge_pipeline.py}; (ii)~ONNX model export and INT8 quantization in \path{pipeline/export_edge_models.py}; (iii)~controlled baselines in \path{experiments/baseline_comparison/run_controlled_baselines.py}; (iv)~canonical PlantVillage validation in \path{experiments/transductive_validation/run_transductive_validation.py}; (v)~routing and gating evaluation in \path{experiments/routing_calibration_validation/run_routing_validation.py}; (vi)~augmentation sensitivity in \path{experiments/augmentation_sensitivity/run_augmentation_sensitivity.py}; and (vii)~mobile profiling reports in \path{experiments/edge_device_benchmark/results/benchmark_report.json}. Edge deployment instructions and pre-quantized models are provided in \path{DEPLOYMENT_GUIDE.md} and \path{models/}.

\subsection{Provenance and Numerical Reproducibility}
Table~\ref{tab:provenance} documents the provenance of reported values (``Re-run'': evaluated locally via reproduction scripts; ``As reported'': transcribed from literature). Controlled reruns reproduce canonical inductive accuracy exactly (95.08\%) and transductive accuracy within 0.04\,pp (97.37\% vs.\ 97.41\% reference, reflecting \texttt{lsqr} solver floating-point sensitivity across library releases, well below the 0.71\,pp standard deviation). Canonical reference values are reported for 11-class benchmarks to ensure consistency with paired significance and sensitivity analyses (Sections~\ref{sec:significance} and~\ref{sec:aug_sensitivity}).

\begin{table}[htbp]
\caption{Provenance of reported values.}
\label{tab:provenance}
\centering

\footnotesize
\begin{tabularx}{\columnwidth}{l X}
\toprule
\textbf{Value} & \textbf{Provenance} \\
\midrule
Table~\ref{tab:comparative} baselines & Re-run, shared splits/seeds; feature branches follow each method \\
Table~\ref{tab:comparative} BioDINO 11-class & Canonical reference run (reproduced $\le$0.04\,pp) \\
Table~\ref{tab:comparative} BioDINO 5-way & Controlled re-run (this work) \\
Table~\ref{tab:shot_ladder} literature rows & As reported, not re-run \\
Table~\ref{tab:plantdoc_context} & Our PlantDoc field-domain run (seeded 8-view ladder, protocol above) \\
Table~\ref{tab:aug_sensitivity} sensitivity & Re-run, this work (same 30-episode protocol) \\
\bottomrule
\end{tabularx}
\end{table}

